\kchapter{Packaging: Helm and Kustomize}{helm}

\begin{objectives}
\begin{itemize}
  \item why plain YAML stops scaling across environments;
  \item Kustomize: bases, overlays, patches, generators;
  \item Helm: charts, values, templates, releases, repositories;
  \item when to choose which (and how to use both).
\end{itemize}
\end{objectives}

\section{The problem}

Your app needs the same manifests in dev, staging, and production, but with
different replica counts, image tags, resource sizes, and hostnames. Copying
folders and editing by hand quickly leads to drift and mistakes. Two tools
dominate:

\begin{center}
\ktablestart
\begin{tabularx}{\linewidth}{P{2.4cm} L L}
\kheader \kth{} & \kth{Kustomize} & \kth{Helm} \\
Idea & Patch plain YAML with overlays. & Render templates with values. \\
Templating & None --- pure YAML. & Go templates + Sprig functions. \\
Built in & Yes: \kc{kubectl apply -k}. & Separate CLI. \\
Distribution & Git folders. & Versioned charts in repos/OCI registries. \\
Release tracking & No. & Yes: install, upgrade, rollback, history. \\
Best for & Your own apps across environments. & Sharing apps; installing third-party software. \\
\end{tabularx}
\end{center}

\section{Kustomize}

\subsection{Layout}

\begin{plaincode}{folder layout}
k8s/
  base/
    deployment.yaml
    service.yaml
    kustomization.yaml
  overlays/
    dev/
      kustomization.yaml
    prod/
      kustomization.yaml
      replicas-patch.yaml
\end{plaincode}

\begin{yamlcode}{k8s/base/kustomization.yaml}
apiVersion: kustomize.config.k8s.io/v1beta1
kind: Kustomization
resources:
  - deployment.yaml
  - service.yaml
labels:
  - pairs:
      app.kubernetes.io/part-of: shop
    includeSelectors: false
configMapGenerator:
  - name: app-config
    literals:
      - LOG_LEVEL=info
\end{yamlcode}

\begin{yamlcode}{k8s/overlays/prod/kustomization.yaml}
apiVersion: kustomize.config.k8s.io/v1beta1
kind: Kustomization
namespace: shop-prod
resources:
  - ../../base
images:
  - name: ghcr.io/example/api          # rewrite the image tag everywhere
    newTag: "1.4.0"
patches:
  - path: replicas-patch.yaml
  - target:
      kind: Deployment
      name: api
    patch: |-
      - op: replace
        path: /spec/template/spec/containers/0/resources/limits/memory
        value: 1Gi
configMapGenerator:
  - name: app-config
    behavior: merge
    literals:
      - LOG_LEVEL=warn
\end{yamlcode}

\begin{yamlcode}{k8s/overlays/prod/replicas-patch.yaml}
apiVersion: apps/v1
kind: Deployment
metadata:
  name: api
spec:
  replicas: 4
\end{yamlcode}

\begin{bashcode}{terminal}
kubectl kustomize k8s/overlays/prod      # render to stdout and inspect
kubectl apply -k k8s/overlays/prod       # render and apply
kubectl diff  -k k8s/overlays/prod
\end{bashcode}

\begin{keyidea}
\kc{configMapGenerator} appends a content hash to the ConfigMap name
(\kc{app-config-7g2k9}) and rewrites references to it. Change a value and the
Deployment's Pod template changes too --- triggering a rolling update
automatically.
\end{keyidea}

\section{Helm}

\kicon{helm}~\textbf{Helm} is the package manager for Kubernetes. Its
vocabulary:

\begin{itemize}
  \item a \kterm{chart} --- a package of templated manifests;
  \item \kterm{values} --- the inputs that customise a chart;
  \item a \kterm{release} --- one installed instance of a chart in a cluster;
  \item a \kterm{repository} --- where charts are published (HTTP or OCI registry).
\end{itemize}

\subsection{Using other people's charts}

\begin{bashcode}{terminal}
helm repo add bitnami https://charts.bitnami.com/bitnami
helm search repo redis
helm show values bitnami/redis > redis-values.yaml     # see every knob

helm install cache bitnami/redis -n cache --create-namespace -f redis-values.yaml
helm list -A
helm upgrade cache bitnami/redis -n cache -f redis-values.yaml --set replica.replicaCount=2
helm history cache -n cache
helm rollback cache 1 -n cache
helm uninstall cache -n cache

# OCI registries work too
helm install eg oci://docker.io/envoyproxy/gateway-helm --version v1.5.0 -n envoy-gateway-system --create-namespace
\end{bashcode}

\subsection{Writing your own chart}

\begin{bashcode}{terminal}
helm create api      # scaffolds a working chart
\end{bashcode}

\begin{plaincode}{chart layout}
api/
  Chart.yaml            # name, version, appVersion
  values.yaml           # default values
  templates/
    _helpers.tpl        # reusable named templates
    configmap.yaml
    deployment.yaml
    service.yaml
    hpa.yaml
    NOTES.txt           # printed after install
\end{plaincode}

\begin{yamlcode}{api/values.yaml}
replicaCount: 2
image:
  repository: ghcr.io/example/api
  tag: ""                  # defaults to Chart.appVersion
  pullPolicy: IfNotPresent
service:
  port: 80
resources:
  requests: { cpu: 100m, memory: 128Mi }
  limits:   { memory: 256Mi }
autoscaling:
  enabled: false
  minReplicas: 2
  maxReplicas: 10
\end{yamlcode}

\begin{plaincode}{api/templates/deployment.yaml}
apiVersion: apps/v1
kind: Deployment
metadata:
  name: {{ include "api.fullname" . }}
  labels:
    {{- include "api.labels" . | nindent 4 }}
spec:
  {{- if not .Values.autoscaling.enabled }}
  replicas: {{ .Values.replicaCount }}
  {{- end }}
  selector:
    matchLabels:
      {{- include "api.selectorLabels" . | nindent 6 }}
  template:
    metadata:
      labels:
        {{- include "api.selectorLabels" . | nindent 8 }}
      annotations:
        checksum/config: {{ include (print $.Template.BasePath "/configmap.yaml") . | sha256sum }}
    spec:
      containers:
        - name: api
          image: "{{ .Values.image.repository }}:{{ .Values.image.tag | default .Chart.AppVersion }}"
          imagePullPolicy: {{ .Values.image.pullPolicy }}
          ports:
            - name: http
              containerPort: 8000
          resources:
            {{- toYaml .Values.resources | nindent 12 }}
\end{plaincode}

\begin{bashcode}{terminal}
helm lint ./api
helm template my-api ./api -f values-prod.yaml      # render locally, no cluster needed
helm install my-api ./api -n shop --create-namespace
helm upgrade --install my-api ./api -n shop -f values-prod.yaml --atomic --wait
helm package ./api && helm push api-0.1.0.tgz oci://ghcr.io/example/charts
\end{bashcode}

\begin{protip}
In CI/CD always use \kc{helm upgrade --install --atomic --wait}: it installs or
upgrades idempotently, waits until resources are ready, and rolls back
automatically on failure.
\end{protip}

\section{Which one should I use?}

A common, pragmatic split:

\begin{itemize}
  \item \textbf{Helm} to install third-party software (Prometheus, cert-manager,
        Argo CD, databases).
  \item \textbf{Kustomize} (or a small Helm chart) for your own services.
  \item Kustomize can even inflate Helm charts (\kc{helmCharts:} field), and
        Argo CD understands both natively.
\end{itemize}

\begin{recap}
\begin{itemize}
  \item Kustomize: base + overlays + patches, template-free, built into kubectl.
  \item Helm: charts + values = releases, with upgrade/rollback history.
  \item Generators/checksums turn config changes into rolling updates.
\end{itemize}
\end{recap}
